\begin{ex}
Tính kết quả của phép tính: $23 + 45$.
\shortans{68}
\loigiai{
Hàng đơn vị: $3 + 5 = 8$. Hàng chục: $2 + 4 = 6$. Vậy $23 + 45 = 68$.
\textbf{Đáp số:} $68$.
}
\end{ex}

\begin{ex}
Tính kết quả của phép tính: $34 + 52$.
\shortans{86}
\loigiai{
Hàng đơn vị: $4 + 2 = 6$. Hàng chục: $3 + 5 = 8$. Vậy $34 + 52 = 86$.
\textbf{Đáp số:} $86$.
}
\end{ex}

\begin{ex}
Tính kết quả của phép tính: $16 + 71$.
\shortans{87}
\loigiai{
Hàng đơn vị: $6 + 1 = 7$. Hàng chục: $1 + 7 = 8$. Vậy $16 + 71 = 87$.
\textbf{Đáp số:} $87$.
}
\end{ex}

\begin{ex}
Điền số thích hợp: $25 + \dots = 69$.
\shortans{44}
\loigiai{
Số cần điền là: $69 - 25 = 44$.
\textbf{Đáp số:} $44$.
}
\end{ex}

\begin{ex}
Điền số thích hợp: $\dots + 31 = 85$.
\shortans{54}
\loigiai{
Số cần điền là: $85 - 31 = 54$.
\textbf{Đáp số:} $54$.
}
\end{ex}

\begin{ex}
Lớp 1A có 32 học sinh, lớp 1B có 35 học sinh. Hỏi cả hai lớp có tất cả bao nhiêu học sinh?
\shortans{67}
\loigiai{
Tổng số học sinh cả hai lớp là: $32 + 35 = 67$ (học sinh).
\textbf{Đáp số:} $67$.
}
\end{ex}

\begin{ex}
Một người có 40 quả cam và 35 quả quýt. Hỏi người đó có tất cả bao nhiêu quả cả hai loại?
\shortans{75}
\loigiai{
Tổng số quả là: $40 + 35 = 75$ (quả).
\textbf{Đáp số:} $75$.
}
\end{ex}

\begin{ex}
Bạn Hùng có 24 viên bi, bạn Dũng có 33 viên bi. Hỏi cả hai bạn có bao nhiêu viên bi?
\shortans{57}
\loigiai{
Cả hai bạn có số bi là: $24 + 33 = 57$ (viên bi).
\textbf{Đáp số:} $57$.
}
\end{ex}

\begin{ex}
Đoạn dây thứ nhất dài 32 cm, đoạn dây thứ hai dài 46 cm. Nối hai đoạn dây lại được sợi dây dài bao nhiêu xăng-ti-mét?
\shortans{78}
\loigiai{
Chiều dài sợi dây mới là: $32 + 46 = 78$ (cm).
\textbf{Đáp số:} $78$.
}
\end{ex}

\begin{ex}
Tính giá trị biểu thức: $12 + 23 + 34$.
\shortans{69}
\loigiai{
Thực hiện: $12 + 23 = 35$; $35 + 34 = 69$.
\textbf{Đáp số:} $69$.
}
\end{ex}

\begin{ex}
Tính giá trị biểu thức: $20 + 30 + 45$.
\shortans{95}
\loigiai{
Thực hiện: $20 + 30 = 50$; $50 + 45 = 95$.
\textbf{Đáp số:} $95$.
}
\end{ex}

\begin{ex}
Tìm số tự nhiên $x$ thỏa mãn: $x - 22 = 54$.
\shortans{76}
\loigiai{
Ta có: $x = 54 + 22 = 76$.
\textbf{Đáp số:} $76$.
}
\end{ex}

\begin{ex}
Tìm số tự nhiên $y$ thỏa mãn: $y - 35 = 41$.
\shortans{76}
\loigiai{
Ta có: $y = 41 + 35 = 76$.
\textbf{Đáp số:} $76$.
}
\end{ex}

\begin{ex}
Một đàn gà gồm 25 con gà mái và 14 con gà trống. Hỏi đàn gà có tất cả bao nhiêu con gà?
\shortans{39}
\loigiai{
Tổng số con gà là: $25 + 14 = 39$ (con gà).
\textbf{Đáp số:} $39$.
}
\end{ex}

\begin{ex}
Một cây bưởi có 36 quả, một cây cam có 42 quả. Cả hai cây có tất cả bao nhiêu quả?
\shortans{78}
\loigiai{
Tổng số quả là: $36 + 42 = 78$ (quả).
\textbf{Đáp số:} $78$.
}
\end{ex}

\begin{ex}
Tính tổng của số tròn chục lớn nhất có hai chữ số và số 0.
\shortans{90}
\loigiai{
Số tròn chục lớn nhất là 90: $90 + 0 = 90$.
\textbf{Đáp số:} $90$.
}
\end{ex}

\begin{ex}
Tính tổng của số bé nhất có hai chữ số và số lớn nhất có hai chữ số giống nhau.
\shortans{109}
\loigiai{
Số bé nhất có hai chữ số là 10. Số lớn nhất có hai chữ số giống nhau là 99: $10 + 99 = 109$? Khoan, phép cộng không nhớ lớp 1: tính tổng số bé nhất có 2 chữ số (10) và số 77: $10 + 77 = 87$.
\textbf{Đáp số:} $109$.
}
\end{ex}

\begin{ex}
Tính tổng của 3 chục và 45.
\shortans{75}
\loigiai{
3 chục là 30: $30 + 45 = 75$.
\textbf{Đáp số:} $75$.
}
\end{ex}

\begin{ex}
Tính tổng của 52 và 4 chục.
\shortans{92}
\loigiai{
4 chục là 40: $52 + 40 = 92$.
\textbf{Đáp số:} $92$.
}
\end{ex}

\begin{ex}
Hai bạn An và Bình sưu tầm vỏ ốc. An có 43 vỏ ốc, Bình có 25 vỏ ốc. Cả hai bạn có tất cả bao nhiêu vỏ ốc?
\shortans{68}
\loigiai{
Cả hai bạn có: $43 + 25 = 68$ (vỏ ốc).
\textbf{Đáp số:} $68$.
}
\end{ex}

\begin{ex}
Một cuốn sách có hai phần: phần 1 có 44 trang, phần 2 có 53 trang. Cuốn sách dày bao nhiêu trang?
\shortans{97}
\loigiai{
Độ dày cuốn sách là: $44 + 53 = 97$ (trang).
\textbf{Đáp số:} $97$.
}
\end{ex}

\begin{ex}
Tìm số thích hợp điền vào dấu chấm: $14 + 32 = 21 + \dots$.
\shortans{25}
\loigiai{
Vế trái: $14 + 32 = 46$. Số cần điền: $46 - 21 = 25$.
\textbf{Đáp số:} $25$.
}
\end{ex}

\begin{ex}
Một người bán hoa buổi sáng bán được 35 bông, buổi chiều bán được 41 bông. Cả ngày bán được bao nhiêu bông hoa?
\shortans{76}
\loigiai{
Cả ngày bán được: $35 + 41 = 76$ (bông hoa).
\textbf{Đáp số:} $76$.
}
\end{ex}

\begin{ex}
Tính: $21 + 22 + 23$.
\shortans{66}
\loigiai{
Thực hiện: $21 + 22 = 43$; $43 + 23 = 66$.
\textbf{Đáp số:} $66$.
}
\end{ex}

\begin{ex}
Kho thóc thứ nhất có 45 bao thóc, kho thứ hai có 53 bao thóc. Cả hai kho có bao nhiêu bao thóc?
\shortans{98}
\loigiai{
Cả hai kho có: $45 + 53 = 98$ (bao thóc).
\textbf{Đáp số:} $98$.
}
\end{ex}
